\section{Reformulating the model as a Hamiltonian system}\label{sec:hamiltonian-and-symmetries}

In this section, we apply the covariant phase-space formalism to the Maxwell theory with the falloffs of Section \ref{sec:maxwell-fields-global-ads3}. We first rewrite the model as a Hamiltonian system and then discuss the Noether charges associated with the Killing symmetries. Appendix \ref{app:falloffs-variational-principle} shows that the near-boundary contributions at spatial infinity vanish.

\subsection{The Hamiltonian system}\label{subsec:hamiltonian-system}

Taking the variation of the Maxwell action gives

\begin{align}
\delta S & =\int _{M}\mathrm{d}^{3}x\sqrt{ -g }\,\nabla_{\nu}F^{\nu \mu}\delta A_{\mu} \\
 & +\int _{\Sigma _{f}}\mathrm{d}^{2}x\sqrt{ \sigma }\,\tau _{\mu}F^{\mu \nu}\delta A_{\nu}-\int _{\Sigma _{i}}\mathrm{d}^{2}x\sqrt{ \sigma }\,\tau _{\mu}F^{\mu \nu}\delta A_{\nu}.
\end{align}

Here $\sigma_{ab}$ is the induced metric on a constant-$t$ Cauchy surface:

\begin{align}
\mathrm{d}s^{2}_{\Sigma} & =\sigma _{ab}\mathrm{d}x^{a}\mathrm{d}x^{b} \\
&=\frac{\mathrm{d}r^{2}}{1+r^{2}}+r^{2}\mathrm{d}\phi ^{2}.
\end{align}

The future-pointing unit normal vector orthogonal to this surface is

\begin{align}
\tau ^{\mu}&=\frac{1}{\sqrt{ 1+r^{2} }}\delta _{0}^{\mu}.
\end{align}

The boundary term produced at spatial infinity vanishes under the asymptotic conditions of Section \ref{sec:maxwell-fields-global-ads3} together with the radial derivative behavior summarized in Appendix \ref{app:falloffs-variational-principle}. From the bulk term we read off the equation of motion

\begin{align}
E^{\mu}&\equiv \nabla _{\nu}F^{\nu \mu}=0.
\end{align}

The symplectic potential on a Cauchy surface $\Sigma$ is

\begin{align}
\theta & =\int _{\Sigma}\mathrm{d}^{2}x\sqrt{ \sigma }\tau _{\mu}F^{\mu \nu}\delta A_{\nu}
\end{align}

We define the pre-phase space $\widetilde{\mathcal{P}}$ as the space of solutions of the Maxwell equation obeying the asymptotic boundary conditions of Section \ref{sec:maxwell-fields-global-ads3}. Taking the exterior derivative on configuration space gives the symplectic form

\begin{align}
\omega & =\delta \theta \\
 & =\int _{\Sigma}\mathrm{d}^{2}x\sqrt{ \sigma }\tau _{\mu}\delta F^{\mu \nu}\wedge \delta A_{\nu}
\end{align}

Appendix \ref{app:falloffs-variational-principle} shows that this symplectic form is finite on $\widetilde{\mathcal{P}}$. At this stage, the Maxwell theory has been rewritten as a Hamiltonian system specified by $\widetilde{\mathcal{P}}$ and $\omega|_{\widetilde{\mathcal{P}}}$.

Because the theory is free, $\widetilde{\mathcal{P}}$ is naturally a linear space. We therefore identify a solution $A_{i,\mu}$ with the corresponding tangent vector on $\widetilde{\mathcal{P}}$,

\begin{align}
X_{A_{i}}&=\int \mathrm{d}^{3}x\,A_{i,\mu}(x)\frac{\delta}{\delta A_{\mu}(x)},
\end{align}

and evaluate the symplectic form on two such tangent vectors to obtain the bilinear pairing

\begin{align}
\Omega[A_{1},A_{2}]&=X_{A_{2}}\cdot X_{A_{1}}\cdot \omega|_{\widetilde{\mathcal{P}}} \\
 & =\int _{\Sigma}\mathrm{d}^{2}x\sqrt{ \sigma }\,\tau _{\mu}\left(F_{1}^{\mu \nu}A_{2,\nu}-F_{2}^{\mu \nu}A_{1,\nu}\right).
\end{align}

This bilinear form is finite and independent of the choice of Cauchy surface. Later we will use $i\Omega[A_{1},A_{2}^{*}]$ as the natural Hermitian inner product on the one-particle space.

\subsection{Noether theorem for Killing symmetries}\label{subsec:noether-killing-symmetries}

We now discuss the Noether charges associated with the Killing symmetries. These charges will later give the Hamiltonian and the angular momentum of the mode expansion.

A Killing symmetry is represented on configuration space by the vector field

\begin{align}
X_{\xi}&=\int \mathrm{d}^{3}x\,\mathcal{L}_{\xi}A_{\mu}\frac{\delta}{\delta A_{\mu}},
\end{align}

where the infinitesimal variation is generated by the Lie derivative,

\begin{align}
X_{\xi}\cdot \delta A_{\mu}&=\mathcal{L}_{\xi}A_{\mu}\equiv \xi ^{\nu}\nabla _{\nu}A_{\mu}+A_{\nu}\nabla _{\mu}\xi ^{\nu}.
\end{align}

Here $\xi^{\mu}$ is a Killing field satisfying

\begin{align}
\nabla _{\mu}\xi _{\nu}+\nabla _{\nu}\xi _{\mu}&=0.
\end{align}

For the global AdS$_3$ Killing fields introduced in Section \ref{sec:maxwell-fields-global-ads3}, one has the asymptotic behavior

\begin{align}
\xi ^{t}&=O(r^{0}), & \xi ^{r}&=O(r), & \xi ^{\phi}&=O(r^{0}),
\end{align}

Acting this symmetry on the action yields

\begin{align}
X_{\xi}\cdot \delta S & =\int _{\Sigma _{f}}\mathrm{d}^{2}x\sqrt{ \sigma }\,\tau _{\rho}\xi ^{\rho}\frac{1}{4}F_{\mu \nu}F^{\mu \nu}-\int _{\Sigma _{i}}\mathrm{d}^{2}x\sqrt{ \sigma }\,\tau _{\rho}\xi ^{\rho}\frac{1}{4}F_{\mu \nu}F^{\mu \nu}.
\end{align}

The contribution from spatial infinity vanishes, so we identify

\begin{align}
\alpha _{\xi}&=\int _{\Sigma}\mathrm{d}^{2}x\sqrt{ \sigma }\,\tau _{\rho}\xi ^{\rho}\frac{1}{4}F_{\mu \nu}F^{\mu \nu}.
\end{align}

Following the covariant phase-space prescription, the corresponding Noether charge on $\widetilde{\mathcal{P}}$ is

\begin{align}
H_{\xi}|_{\widetilde{\mathcal{P}}}&=(X_{\xi}\cdot \theta-\alpha _{\xi})|_{\widetilde{\mathcal{P}}} \\
 & =\left.\int _{\Sigma}\mathrm{d}^{2}x\sqrt{ \sigma }\,\tau _{\mu}\left(F^{\mu \nu}\mathcal{L}_{\xi}A_{\nu}-\frac{1}{4}\xi ^{\mu}F_{\rho \sigma}F^{\rho \sigma}\right)\right|_{\widetilde{\mathcal{P}}} \\
 & =\left.\int _{\Sigma}\mathrm{d}^{2}x\sqrt{ \sigma }\,\tau _{\mu}\xi _{\nu}\left(F^{\mu \rho}F^{\nu}{}_{\rho}-\frac{1}{4}g^{\mu \nu}F_{\rho \sigma}F^{\rho \sigma}\right)\right|_{\widetilde{\mathcal{P}}} \\
 & +\left.\int _{\Sigma}\mathrm{d}^{2}x\sqrt{ \sigma }\,\tau _{\mu}\left[\nabla _{\nu}\left(F^{\mu \nu}A_{\rho}\xi ^{\rho}\right)+\left(\nabla _{\nu}F^{\nu \mu}\right)A_{\rho}\xi ^{\rho}\right]\right|_{\widetilde{\mathcal{P}}} \\
 & =\left.\int _{\Sigma}\mathrm{d}^{2}x\sqrt{ \sigma }\,\tau _{\mu}\xi _{\nu}\left(F^{\mu \rho}F^{\nu}{}_{\rho}-\frac{1}{4}g^{\mu \nu}F_{\rho \sigma}F^{\rho \sigma}\right)\right|_{\widetilde{\mathcal{P}}}.
\end{align}

In the last step we use the Maxwell equation and the vanishing of the total-derivative term at the boundary of $\Sigma$.

We next verify the Hamiltonian relation between the Killing flow, the symplectic form, and the Noether charge. Using the Lie derivative on $A_{\mu}$ and on $F^{\mu\nu}$, one finds

\begin{align}
X_{\xi}\cdot \omega|_{\tilde{\mathcal{P}}}&=\int _{\Sigma}\mathrm{d}^{2}x\sqrt{ \sigma }\,\tau _{\mu}\left((\mathcal{L}_{\xi}F^{\mu \nu})\delta A_{\nu}-(\mathcal{L}_{\xi}A_{\nu})\delta F^{\mu \nu}\right)\bigg|_{\tilde{\mathcal{P}}} \\
 & =-\delta H_{\xi}|_{\tilde{\mathcal{P}}}
\end{align}

For later use, it is convenient to record three immediate consequences. First, for two linearized solutions $A_{1}$ and $A_{2}$ one has

\begin{align}
\Omega[\mathcal{L}_{\xi}A_{1},A_{2}]&=-\Omega[A_{1},\mathcal{L}_{\xi}A_{2}]=-H_{\xi}^{(2)}[A_{1},A_{2}],
\end{align}

where

\begin{align}
H_{\xi}^{(2)}[A_{1},A_{2}]&=\int _{\Sigma}\mathrm{d}^{2}x\sqrt{ \sigma }\,\tau _{\mu}\xi _{\nu}\left(F_{1}^{\mu \rho}F_{2,\rho}^{\nu}+F_{2}^{\mu \rho}F_{1,\rho}^{\nu}-\frac{1}{2}g^{\mu \nu}F_{1,\rho \sigma}F_{2}^{\rho \sigma}\right).
\end{align}

This bilinear form satisfies $H_{\xi}^{(2)}[A,A]=2H_{\xi}[A]$. Second, the Noether charge itself may be written as

\begin{align}
H_{\xi}[A]&=-\frac{1}{2}\Omega[\mathcal{L}_{\xi}A,A]=\frac{1}{2}\Omega[A,\mathcal{L}_{\xi}A].
\end{align}

Third, the symplectic pairing is invariant under the Killing flow:

\begin{align}
\Omega[\mathcal{L}_{\xi}A_{1},A_{2}]+\Omega[A_{1},\mathcal{L}_{\xi}A_{2}]&=0.
\end{align}

These identities provide the bridge between the Hamiltonian formulation of this section and the mode expansion carried out in Sections 3 and 4.
